% From mitthesis package
% Version: 1.01, 2023/07/04
% Documentation: https://ctan.org/pkg/mitthesis


\chapter{Code listing}

% This example uses the \texttt{listings} package.

\bigskip

\lstdefinestyle{mystyle}{
    backgroundcolor=\color{gray!5},   
    commentstyle=\color{gray}\itshape,
    numberstyle=\tiny\color{gray!50},
    stringstyle=\color{black},
    basicstyle=\small\ttfamily,
    breakatwhitespace=false,         
    breaklines=true,                 
    numbers=left,                    
    numbersep=5pt,                  
    showspaces=false,                
    showstringspaces=false,
    showtabs=false,                  
    tabsize=2
}%

\lstset{language=Matlab, style={mystyle}}%
\subsubsection{Functions used for processing behavioural data}
\begin{lstlisting}
    function behavProcess(myFolder, pattern, destFin, timeprd, incmt, rescaleData, deadzone, normalizeData, dzFac, val_ind, arl_ind, RangeMax, RangeMin, InputMax, InputMin)
    % The following code takes in raw participant behavioural data from
    % specified location (myFolder), downsamples it to specified timeperiod
    % (timeprd), with specified increments(incmts). Then it rescales the
    % downsampled participant data to a specified range (RangeMax,
    % RangeMin) from the inherent input range (InputMax, InputMin). If the
    % deadzone variable is true it then applies a deadzone at the center of
    % a specified threshold (dzFac). Finally it writes the processed data
    % to a file named sub-%%_processed.csv

    % Read in the file location details
    filePattern = fullfile(myFolder, pattern); % Change to whatever pattern you need.
    theFiles = dir(filePattern);

    % Initialization of variables
    timescale = (0:incmt:timeprd)';
    val_data = zeros(length(timescale),length(theFiles));
    arl_data = zeros(length(timescale),length(theFiles));
    normVal = zeros(length(timescale),length(theFiles));
    normArl = zeros(length(timescale),length(theFiles));
    valData = zeros(length(timescale),length(theFiles));
    arlData = zeros(length(timescale),length(theFiles));
    
    % DOWNSAMPLING
    % Loop over file list and read in the behaviour data
    for k = 1 : length(theFiles)
        baseFileName = theFiles(k).name;
        fullFileName = fullfile(theFiles(k).folder, baseFileName);
        fprintf(1, 'Now reading %s\n', baseFileName);

        % Read in the file
        emo_data = readmatrix(fullFileName);

        %Loop over the num of observations in the final table
        
        for i = 1:length(timescale)
            % Timescale is a incmt multiple of every observation num
            cond = round((i-1)*incmt,2);
            % Check which timepoints come close to the established timepoint (cond)
            [minVal,minInd] = min(abs(emo_data(:,2)-cond));
            % Index out closest timepoint from the original data
            val_data(i,k) = emo_data(minInd,val_ind);
            arl_data(i,k) = emo_data(minInd,arl_ind);
        end
    end
    
    % RESCALING
    % Rescale the valence and arousal datapoints to specified range
    if rescaleData == true
        for k = 1 : length(theFiles)
            % Rescales the data to specified range
            valData(:,k) = rescale(val_data(:,k),RangeMin,RangeMax,"InputMax",InputMax,"InputMin",InputMin);
            arlData(:,k) = rescale(arl_data(:,k),RangeMin,RangeMax,'InputMax',InputMax,'InputMin',InputMin);
        end
    else
        valData = val_data;
        arlData = arl_data;
    end
    
    % DEADZONING
    % Apply deadzone and center participant data across participants
    % within time points
    if deadzone == true
        for k = 1 : length(theFiles)
            for rowInd = 1:length(timescale)
                % Anything between +/- deadzone filter is made zero for
                % both valence and arousal
                if valData(rowInd,k) <= dzFac && valData(rowInd,k) >= -dzFac
                   valData(rowInd,k) = 0;
                end
                if arlData(rowInd,k) <= dzFac && arlData(rowInd,k) >= -dzFac
                   arlData(rowInd,k) = 0;
                end
            end
        end
    end
    
    % NORMALIZATION
    if normalizeData == true
        % Normalize participant responses
        valData = normalize(valData, 1,"zscore");
        arlData = normalize(arlData, 1,"zscore");
    end
    
    % SAVING FINAL DATA
    % Output the final data and save it to specified file format
    for k = 1 : length(theFiles)
        baseFileName = theFiles(k).name;
        name = split(baseFileName,"_");
        fprintf(1, 'Now preparing to write processed data for subject %s\n', name{1});        
        final_table = zeros(length(timescale),3);
        final_table(:,1) = timescale; % Timescale created by us
        final_table(:,2) = valData(:,k); % Initial 5 time points are zero
        final_table(:,3) = arlData(:,k); % Initial 5 time points are zero
        %final_table(1,:) = zeros(1,width(final_table));
        resultname = strcat(destFin, name{1},'_processed.csv');
        writematrix(final_table,resultname,"WriteMode","overwrite");
    end
end


function calcKinematics(myFolder,pattern, val_ind, arl_in,time_ind, strength_ind, disp_ind, vel_ind, acc_ind, jerk_ind, filter, fltVal)
% The function takes in cleaned and pre-processed behaviour time series and
% calculates the kinematics for that participant

% myFolder specifies the path to be checked for sources files, pattern
% tells the function the name pattern of the files to be looked for
% val_ind, arl_in, time_ind specify the indices for valence, arousal and
% the time in the source files. Displacement, velocity and acceleration
% indices specify the indices where the corresponding kinematics values are
% to filled in. 

% Filter specifies whether the kinematics should be thresholded on the
% basis of displacement distribution for the participant. fltVal specifies
% the number of standard deviations the kinematics distribution should be
% filtered within.

    % Get a list of all files in the folder with the desired file name pattern.
    filePattern = fullfile(myFolder, pattern); % Change to whatever pattern you need.
    theFiles = dir(filePattern);
    
    % Loop through the file list
    for k = 1 : length(theFiles)
        baseFileName = theFiles(k).name;
        fullFileName = fullfile(theFiles(k).folder, baseFileName);
        fprintf(1, 'Now reading %s\n', fullFileName);
    
        % Read in the participant data
        emo_data = readmatrix(fullFileName);
        
        % Initialize empty rows and fill them with zeros, to be filled later
        emo_data(1,disp_ind) = zeros();% Displacement
        emo_data(1,vel_ind) = zeros();% Velocity
        emo_data(1,acc_ind) = zeros();% Acceleration
        
        fprintf(1, 'Calculating resultant displacement, velocity and accelaration \n')
        
        % Indexes from 2 because Intensity at t = 0/1 will be 0
        for i = 1:length(emo_data)
            % Calculate the intensity of emotion at each
            % consecutive time frame.
            % Intensity  = sqrt (valence(t)^2 + arousal(t)^2)
            emo_data(i,strength_ind) = sqrt( (emo_data(i,val_ind) + emo_data(i,arl_in) )^2 );
        end
        
        % Indexes from 2 because Displacement at t = 0/1 will be 0
        for i = 2:length(emo_data)-1
            % Calculate the displacement of the joystick at each
            % consecutive time frame.
            % Displacement = sqrt( (valence(t) - valence(t-1))^2 + (arousal(t) - arousal(t-1))^2)
            emo_data(i, disp_ind) = sqrt( (emo_data(i+1, val_ind) - emo_data(i-1, val_ind) )^2 +  (emo_data(i+1, arl_in) - emo_data(i-1, arl_in) )^2 );
        end

        if filter == true
            % Calculate the standard deviation of the displacement values
            % for the participant. This will be used as a filter and only
            % displacements "above" fltVal*SD will be allowed to be used
            % for velocity calculation.
            sdThresh = fltVal * std(emo_data(:,disp_ind));
            % Indexes from 2 because velocity at t = 0/1 will be 0
            for i = 3:length(emo_data)-1
                if emo_data(i,disp_ind) <= sdThresh && emo_data(i,disp_ind) >= -sdThresh
                    % Make any timepoint with displacement values less than
                    % sdThresh zero
                    emo_data(i,disp_ind) = 0;
                    %emo_data(i,strength_ind) = 0;
                else
                    % Calculate the velocity of the joystick at each
                    % consecutive time frame.
                    % Velocity = displacement(t-1:t)/time(t-(t-1))
                    emo_data(i, vel_ind) = emo_data(i, disp_ind) / 2*((emo_data(i, time_ind) - emo_data(i-1, time_ind)));
                end
            end
        else
            for i = 3:length(emo_data)-1
            % If filter is not true velocity is calculated unthresholded
            % Velocity = displacement(t-1:t)/time(t-(t-1))
                emo_data(i, vel_ind) = emo_data(i, disp_ind) / 2*((emo_data(i, time_ind) - emo_data(i-1, time_ind)));
            end
        end
        
        % Indexes from 4 because Acceleration at initial two time points will be 0
        for i = 4:length(emo_data)-1
            % Calculate the acceleration 
            % Acceleration = (velocity(t) - velocity(t-1))/time(t-(t-1))
            emo_data(i,acc_ind) = (emo_data(i+1, vel_ind) - emo_data(i-1, vel_ind)) / 2*(emo_data(i, time_ind) - emo_data(i-1, time_ind));
        end
        

        for i = 5:length(emo_data)-1
            % Calculate the acceleration 
            % Acceleration = (velocity(t) - velocity(t-1))/time(t-(t-1))
            emo_data(i,jerk_ind) = (emo_data(i+1, acc_ind)-emo_data(i-1, acc_ind)) / 2*(emo_data(i,time_ind)-emo_data(i-1,time_ind));
        end
        % Write the values to the specified file
        writematrix(emo_data, fullFileName,'WriteMode','overwrite');
    
    end
end

function [valence, arousal, strength, displacement, time, velocity, acceleration, jerk] = disFeat_Exact(myFolder,pattern,val_ind, arl_ind, strength_ind, disp_ind, time_ind, vel_ind,acc_ind, jerk_ind)  

% This function has been create to extract state and kinematics information
% from participant files with ease without having to call out each
% individual participant repeatedly. 

    % Get a list of all files in the folder with the desired file name pattern.
    filePattern = fullfile(myFolder, pattern); % Change to whatever pattern you need.
    % Contains the names/locations of files with specified name pattern
    theFiles = dir(filePattern);
    
    % Initialise empty vectors of required length
    valence = [1,length(theFiles)];
    arousal = [1,length(theFiles)];
    strength = [1,length(theFiles)];
    displacement = [1,length(theFiles)];
    velocity = [1,length(theFiles)];
    time = [1,length(theFiles)];
    acceleration = [1,length(theFiles)];
    jerk = [1,length(theFiles)];

    % Loop across files/participant numbers
    for k = 1 : length(theFiles)
        baseFileName = theFiles(k).name;
        fullFileName = fullfile(theFiles(k).folder, baseFileName);
        fprintf(1, 'Now reading %s\n', fullFileName);
        
        % Load participant data files
        emo_fdata = readmatrix(fullFileName);
        
        % Look across each time points
        for i = 1: length(emo_fdata)
            % Indexing out each variable and filling empty arrays
            valence(i,k) = emo_fdata(i,val_ind);
            arousal(i,k) = emo_fdata(i,arl_ind);
            strength(i,k) = emo_fdata(i,strength_ind);
            displacement(i,k) = emo_fdata(i,disp_ind);
            time(i,k) = emo_fdata(i,time_ind);
            velocity(i,k) = emo_fdata(i,vel_ind);
            acceleration(i,k) = emo_fdata(i,acc_ind);
            jerk(i,k) = emo_fdata(i,jerk_ind);
        end
            
    end
    
end


function [actualBoundary, actualFrame] = timepointCalc(timepoints,timeGap, timeStop)
    
    % This is the timescale followed by the data   
    timeFrame = 0:timeGap:timeStop;
    
    % Find the timepoints of shift
    % ragaTime  = [45, 125, 157, 192, 258, 366, 442];
        % This is to calculate the timepoints corresponding to the Raga shift
        % time points

    actualBoundary = zeros(length(timepoints)+2,1);
    for i = 1:length(timepoints)
        [minVal,minInd] = min(abs(timeFrame - timepoints(i)));
        actualBoundary(i+1) = timeFrame(minInd);
    end
    actualBoundary(length(timepoints)+2) = timeFrame(length(timeFrame));
    
    % Find the frame points of shift
        % This is to calculate the frames corresponding to the Raga shift
        % time points
    actualFrame= ones(length(timepoints)+2,1);
    for i = 1:length(timepoints)
        [minVal,minInd] = min(abs(timeFrame - timepoints(i)));
        actualFrame(i+1) = minInd;
    end
    actualFrame(length(timepoints)+2) = length(timeFrame);

end

function accBinary = binarizeAcc(acceleration, fillWindow, addFrame)
% Binarising acceleration time points on basis of zero vs some value
    % fillWindow will fill any gaps in the acceleration time series that
    % occur for less than 1TR with the previous marker

    % addFrame is used to add a 1TR preceding window to all acceleration
    % events. This was to account for any decision making that would go
    % into the acceleration events

    % Initialise the accBinary variable to store binarised acceleration data
    accBinary = zeros(length(acceleration),width(acceleration));
    % Loop across participants
    for j = 1:width(acceleration)
        % Loop across timepoints
        for i = 1:length(acceleration)
            % Convert all non-zero events into acceleration windows (1)
            if acceleration(i,j) ~= 0
                accBinary(i,j) = 1;
            else
            % Convert all zero events into non-acceleration windows (0)
                accBinary(i,j) = 0;
            end
        end
    end
    
    if fillWindow == true
        % Single occurence windows are converted into preceding event
        for j = 1:width(acceleration)
            for i = 2:length(acceleration)-1
                if accBinary(i-1,j)==accBinary(i+1,j)
                    if accBinary(i-1,j)~=accBinary(i,j)
                        if accBinary(i,j) == 0
                            accBinary(i,j) = 1;
                        else
                            accBinary(i,j) = 0;
                        end
                    end
                end
            end
        end
    end
    
    if addFrame == true
        % Considering that the decision for acceleration precedes actual
        % acceleration, a preceding window is ideal to look into underlying
        % computation.
        
        % Adding a 1 TR window before every acceleration time point
        for j = 1:width(acceleration)
            for i = 2:length(acceleration)
                if accBinary(i-1,j)~=accBinary(i,j)
                   if accBinary(i,j)==1
                       accBinary(i-1,j)=1;
                   end
                end
            end
        end
    end

end

function ind_lat = calcLatency(accBinary, actualFrame, timeBin, divide, duration, limit)

    % Initialise empty array
    ind_lat = zeros(8,width(accBinary));
    
    % List of raga indices
    t_ind = actualFrame;
    % t_ind = [1,39,105,132,161,216,306,369];

    % Initialise loop across participants
    for i = 1:width(accBinary)
    
        % Loop to index out raga indices and generate raga specific tiem frames
        for k = 1:8
        
            % For the first index there isn't a need to account for runovers
            if k==1
                for j = t_ind(k):(t_ind(k+1)-1)
                    % Check if the time point is a transition point or not
                    if accBinary(j,i) == 1
                        if divide == true
                            ind_lat(k,i) = (((j-t_ind(k))+1)*timeBin)/duration(k);
                            if limit == true
                                if ind_lat(k,i) > 32 
                                    ind_lat(k,i) = 0; 
                                end
                            end
                        else
                            ind_lat(k,i) = ((j-t_ind(k))+1)*timeBin;
                            if limit == true
                                if ind_lat(k,i) > 32 
                                    ind_lat(k,i) = 0; 
                                end
                            end
                        end
                        break
                    end
                end
            
            % Loop from the second last raga index to the end
            elseif k==8
                    % Chcek if the 8th index timepoint and timepoint+1 is a transition point
                    if accBinary(t_ind(k),i)==1 || accBinary(t_ind(k)+1,i)==1
                    % Find the next closest which is a non-transition index
                        % If Kth index is a transition point
                        if accBinary(t_ind(k),i)==1
                            for j = t_ind(k):length(accBinary)
                                % The closest non-transition index becomes
                                % the new starting point.
                                if accBinary(j,i) == 0
                                    newInd = j;
                                    break
                                end
                            end
                        % If (K+1)th index is a transition point
                        elseif accBinary(t_ind(k)+1,i)==1
                            for j = t_ind(k)+1:length(accBinary)
                                % The closest non-transition index becomes
                                % the new starting point.
                                if accBinary(j,i) == 0
                                    newInd = j;
                                    break
                                end
                            end
                        end
                        
                        % Search from this new index for the closest
                        % transition point. 
                        for j = newInd:length(accBinary)
                            if accBinary(j,i)==1
                                if divide == true
                                    ind_lat(k,i) = (((j-t_ind(k))+1)*timeBin)/duration(k);
                                    if limit == true
                                        if ind_lat(k,i) > 32 
                                            ind_lat(k,i) = 0; 
                                        end
                                    end

                                else
                                    ind_lat(k,i) = ((j-t_ind(k))+1)*timeBin;
                                    if limit == true
                                        if ind_lat(k,i) > 32 
                                            ind_lat(k,i) = 0; 
                                        end
                                    end
                                end
                                break
                            end
                        end
                    
                    % If the Kth and (K+1)th index are not transition
                    % points, directly move to calculating latency.
                    else 
                        for j = t_ind(k):length(accBinary)
                            if accBinary(j,i) == 1
                                if divide == true
                                    ind_lat(k,i) = (((j-t_ind(k))+1)*timeBin)/duration(k);
                                    if limit == true
                                        if ind_lat(k,i) > 32 
                                            ind_lat(k,i) = 0; 
                                        end
                                    end
                                else
                                    ind_lat(k,i) = ((j-t_ind(k))+1)*timeBin;
                                    if limit == true
                                        if ind_lat(k,i) > 32 
                                            ind_lat(k,i) = 0; 
                                        end
                                    end
                                end
                                break
                            end
                        end
                    end
    
            % For other indices
            else
                % Chcek if the 8th index timepoint and timepoint+1 is a transition point
                if accBinary(t_ind(k),i)==1 || accBinary(t_ind(k)+1,i)==1
                % Find the next closest which is a non-transition index
                    % If Kth index is a transition point
                    if accBinary(t_ind(k),i)==1
                        for j = t_ind(k):t_ind(k+1)
                            % The closest non-transition index becomes
                            % the new starting point.
                            if accBinary(j,i) == 0
                                newInd = j;
                                break
                            end
                        end
                    % If (K+1)th index is a transition point    
                    elseif accBinary(t_ind(k)+1,i)==1
                        for j = t_ind(k)+1:t_ind(k+1)
                            % The closest non-transition index becomes
                            % the new starting point.
                            if accBinary(j,i) == 0
                                newInd = j;
                                break
                            end
                        end
                    end

                    % Search from this new index for the closest
                    % transition point. 
                    for j = newInd:t_ind(k+1)
                        if accBinary(j,i)==1
                            if divide == true
                                ind_lat(k,i) = (((j-t_ind(k))+1)*timeBin)/duration(k);
                                if limit == true
                                    if ind_lat(k,i) > 32 
                                        ind_lat(k,i) = 0; 
                                    end
                                end
                            else
                                ind_lat(k,i) = ((j-t_ind(k))+1)*timeBin;
                                if limit == true
                                    if ind_lat(k,i) > 32 
                                        ind_lat(k,i) = 0; 
                                    end
                                end
                            end
                            break
                        end
                    end

                % If the Kth and (K+1)th index are not transition
                % points, directly move to calculating latency.
                else 
                    for j = t_ind(k):t_ind(k+1)
                        if accBinary(j,i) == 1
                            if divide == true
                                ind_lat(k,i) = (((j-t_ind(k))+1)*timeBin)/duration(k);
                                if limit == true
                                    if ind_lat(k,i) > 32 
                                        ind_lat(k,i) = 0; 
                                    end
                                end

                            else
                                ind_lat(k,i) = ((j-t_ind(k))+1)*timeBin;
                                if limit == true
                                    if ind_lat(k,i) > 32 
                                        ind_lat(k,i) = 0; 
                                    end
                                end
                            end
                            break
                        end
                    end
                end
            end
        end
    end
end

\end{lstlisting}


\lstset{language=Python, style={mystyle}}%

\subsubsection{Denoising pre-processed neural data}

\begin{lstlisting}
    #%% Import required packages

from os.path import  join, basename
import glob
import pandas as pd
#import numpy as np
import nilearn.image as image
import nibabel as nib
from nilearn.interfaces.fmriprep import load_confounds

#%% Slicing BOLD files

''' 
    Removal of the first 13 and last 2 volumes from the raw BOLD data (443 volumes) 
    to make it a 428 volume time series.
'''

# Location of EPI files:
data_dir = '/media/austin/Emotion/InterAffectTravel/neuralData/workingBOLD/epi_raw/'

# glob over all possible files with given name pattern in the provided directory:
file_list = sorted(glob.glob(join(data_dir, '*_task-music_space-MNI152NLin2009cAsym_res-2_desc-preproc_bold.nii.gz')))

# Loop through all files in the file list acquired by globbing over given directory:
for f in file_list:
    # Extract subject name:
    sub = basename(f).split('_')[0]
    print(sub)
    # Load the original EPI file
    f_img = image.load_img(f)
    # Index out required volumes and create a new volume
    final_img = image.index_img(f_img,slice(12,440))
    # Save new EPI file
    final_img.to_filename(f)

#%% Slicing the confound files

''' 
    Removal of the first 13 and last 2 volumes from the raw confounds data (443 volumes) 
    to make it a 428 volume time series.
'''

# Location of confound files:
data_dir = '/media/austin/Emotion/InterAffectTravel/neuralData/workingBOLD/epi_raw/'

# glob over all possible files with given name pattern in the provided directory:
file_list = sorted(glob.glob(join(data_dir,'*_task-music_desc-confounds_timeseries.tsv')))

# Loop through all files in the file list acquired by globbing over given directory:
for f in file_list:   
    # Extract subject name:
    sub = basename(f).split('_')[0]
    print(sub)
    # Load the original confound file
    bold_data = pd.read_csv(f, sep='\t', header = 0)
    # Index out required enteries and create a new matrix
    func_data = bold_data[12:440]
    # Save new confound file
    func_data.to_csv(f'/media/austin/Emotion/InterAffectTravel/neuralData/workingBOLD/epi_raw/{sub}_task-music_desc-confounds_timeseries.tsv', index=False, sep='\t')
    
#%% Extracting framewise displacement values and creating a matrix

# Location of confound files:
data_dir = '/media/austin/Emotion/InterAffectTravel/neuralData/workingBOLD/epi_raw/'

# glob over all possible files with given name pattern in the provided directory:
file_list = sorted(glob.glob(join(data_dir,'*_task-music_desc-confounds_timeseries.tsv')))

# Loop through all files in the file list acquired by globbing over given directory:
fd = []
for f in file_list:
    # Extract subject name:
    sub = basename(f).split('_')[0]
    print(sub)
    # Load the original confound file
    bold_data = pd.read_csv(f, sep='\t', header = 0)
    # EXtract out the FD columns
    func_data = bold_data.loc[:,'framewise_displacement']
    # Index out required enteries and create a new matrix
    fd.append(func_data)
# Transpose the created matrix into (timepoints,subject):
fd = pd.DataFrame(fd).T
# Rename each column for participant:
fd = fd.set_axis([f'FD_sub-{x:02}' for x in range(1,len(file_list)+1)], axis=1)
# Save the FD data matrix
fd.to_csv('/media/austin/Emotion/InterAffectTravel/neuralData/workingBOLD/fd_allpart.csv', index=False)

#%% Denoising the data, removing the effects of motion variables, CSF, and detrending the data

''' 
    This script is for loading in the confound files generated by fmriprep during preprocessing and regressing them out
    to denoise the BOLD data. It regresses out CSF effects, WM activity, motion artifacts captured by Framewise displacement and drift regressors from cosine. 
    
'''

# Smoothing kernel
fwhm = 6 

# Location of framewise displacement data matrix, to be used in confound removal
fd_val = pd.read_csv('/media/austin/Emotion/InterAffectTravel/neuralData/workingBOLD/fd_allpart.csv')

# Directory with preprocessed EPI files
base_dir = '/media/austin/Emotion/InterAffectTravel/neuralData/workingBOLD/' 

# globing over all in specified format in the base_dir
file_list = sorted([x for x in glob.glob(join(base_dir, '*epi_raw/*_task-music_space-MNI152NLin2009cAsym_res-2_desc-preproc_bold.nii.gz')) if 'denoise' not in x])

# Looping over each file found by glob
for i, f in enumerate(file_list):    
    
    # Extract subject name
    sub = basename(f).split('_')[0]
    print('The following subject data is being denoised:',sub)
    
    # Load EPI file
    readBrain = nib.load(f)
    
    # Smoothen the image according to specified smoothing kernel
    brain_data = image.smooth_img(readBrain, fwhm = fwhm)
        
    # Load in required confounds, generate a design matrix to regress them out
    confounds_simple, sample_mask = load_confounds(
        f,
        strategy=["motion", "wm_csf", "high_pass"],
        motion="derivatives",
        wm_csf="derivatives",
        demean = True # Centering the confounds for each participant, Made false
    )
    
    # Include the framewise displacement data in the confound matrix
    confound_final = pd.concat([confounds_simple,fd_val.iloc[:,i]], axis = 1)
    print("The shape of the confound matrix is:", confound_final.shape)
    print(confound_final.columns)
        
    # final_img = clean_img(brain_data, detrend=True, confound = confounds_simple, t_r = 1.2, mask_img = brain_mask)
    # Don't load masks ahead of time, nilern GLM instance errors out as it can't figure out where to derive information from
    # Detrending removes the influence of proximate voxels on each other
    # Regress out each variable in the confounds matrix obtaining the final denoised BOLD data:
    final_img = image.clean_img(brain_data, detrend = True, confound = confound_final, t_r = 1.2) 
    
    # Save the final processed BOLD file:
    print("Saving denoised data for subject:", sub)
    nib.save(final_img, join(base_dir, 'epi_CosineRemoved', f'{sub}_nilearn_denoise_smooth{fwhm}mm_task-music_space-MNI152NLin2009cAsym_res-2_desc-preproc_bold.nii.gz'))
\end{lstlisting}


\subsubsection{GLM analysis of the neural data}

\begin{lstlisting}
#%% # Creating the acceleration design matrix revamped

def create_event_file(timeseries, tr=1.0):
    """
    Create event file for Nilearn from a binary timeseries.
    
    Parameters:
    - timeseries: 1D array-like of 0s and 1s
    - tr: Repetition time (in seconds)

    Returns:
    - events: DataFrame with onset, duration, trial_type
    """
    timeseries = np.asarray(timeseries)
    assert set(np.unique(timeseries)).issubset({0, 1}), "Timeseries must only contain 0 and 1"

    events = []
    current_val = timeseries[0]
    onset = 0

    for i in range(1, len(timeseries)):
        if timeseries[i] != current_val:
            duration = round((i - onset) * tr,1) - tr
            trial_type = ["NA" if current_val==0 else "A"][0]
            if duration > 0:
                events.append({
                    "onset": round(onset * tr, 1),
                    "duration": duration,
                    "trial_type": trial_type,
                    "modulation" : 1
                })
            onset = i
            current_val = timeseries[i]

    # Add final event
    duration = round((len(timeseries) - onset) * tr, 1) - tr
    trial_type = ["NA" if current_val==0 else "A"][0]
    if duration > 0:
        events.append({
            "onset": round(onset * tr,1),
            "duration": duration,
            "trial_type": trial_type,
            "modulation" : 1
        })

    return pd.DataFrame(events)

avna = pd.read_csv('/media/austin/Emotion/InterAffectTravel/behaviouralData/Data_TRtime/accBinarised_TRt.csv',header=None).T

conditionS = {}
dm_dict = {}

for sub_ind in range(avna.shape[1]):
    
    conditionS[sub_ind] = create_event_file(avna.iloc[:,sub_ind], tr=1.2)
    
    # Time frame of the whole experiment in seconds, obtained as number of EPI volumes x TR time
    frame_times = np.array([round(x*1.2,2) for x in range(428)])
    # Create the first level design matrix. 
    ''' Glover HRF model is used for our purposes. This can be changed for better fit ''' 
    dm = make_first_level_design_matrix(frame_times, events = conditionS[sub_ind], hrf_model="glover", drift_model=None)
    print(sub_ind)
    # Load generated design matrix to the dictionary with the participant ID as a dictionary key
    dm_dict[sub_ind] = dm

# Save design matrix for figure
plot_design_matrix(dm_dict[3])
plt.savefig("/home/austin/Documents/AustinV_Thesis/figures/finalFigure/GLM_TransitionCoding_DesignMatrix_Sub3.png", transparent=True)


#%% # First Level Analysis GLM run per participant at voxel

# Specify the size of the smoothign kernel used for denoising. 
# THIS IS ONLY USED TO CORRECTLY READ IN FILES
fwhm=6
  
# Specify the directory to glob for denoised EPI files
base_dir = '/media/austin/Emotion/InterAffectTravel/neuralData/workingBOLD/epi_CosineRemoved/'

# glob over provided directory checking for specified condition
file_list = sorted([x for x in glob.glob(os.path.join(base_dir,f'*nilearn_denoise_smooth{fwhm}mm_task-music_space-MNI152NLin2009cAsym_res-2_desc-preproc_bold.nii.gz'))])

# Initiliaze a dictionary to store 1st lvl GLM results
sub_dict = {}

# Loop over file list generated by glob 
for i,f in enumerate(file_list):    
    
    # Extract subject ID
    sub = os.path.basename(f).split('_')[0]
    print('The following subject data is being processed:',sub)
    
    # Load in the EPI file
    neuroData = nib.load(f)
    
    # Call subect specific design matrix, from the dictionary of design matrices created above
    dm = dm_dict[i]
    
    ''' For use only if a whole brain mask template isn't available, otherwise stick to using established templates. Check templateflow if you can't find something specific'''
    # compute a whole brain mask, voxels within this mask will be denoised and extracted.
    #brain_mask = compute_brain_mask(neuroData, mask_type='whole-brain', threshold = 0.08)
    
    ''' Same as above '''
    # Using an MNI152 brain mask to potentially remove the error caused due to different sizes of nifti files and mask
    #brain_mask = mni_mask(resolution=2) #still throwing error
    
    # Load in brain mask template
    brain_mask = nib.load('/media/austin/Emotion/InterAffectTravel/atlases/tpl-MNI152NLin2009cAsym_res-02_desc-brain_mask.nii.gz')
    
    # Initiliaze a first level model
    glm_model = FirstLevelModel(t_r = 1.2, hrf_model='glover', mask_img=brain_mask, n_jobs=-1, signal_scaling=False)
    # Fit brain data to loaded design matrix
    glm_model = glm_model.fit(neuroData, design_matrices = dm)
    
    print("Creating a constrast matrix")
    # Intialize an empty contrast array
    contrast_matrix = np.eye(dm.shape[1])
    
    # Draw contrast values from the design matrix. 
    # Column names and values are drawn and autofilled.
    basic_constrats = {
        column: contrast_matrix[i]
        for i, column in enumerate(dm.columns)
        }
    
    # Setup contrast conditions
    contrasts = {
        "A-NA": basic_constrats["A"] - basic_constrats["NA"],
        "NA-A": basic_constrats["NA"] - basic_constrats["A"],
        "Both": np.vstack((basic_constrats["A"], basic_constrats["NA"]))
        }

    print(f"Computing constrasts for subject: {sub}")
    
    # Compute the contrasts
    for contrast_id, contrast_val in contrasts.items():
        print(f"\tcontrast id: {contrast_id}") 
        #z_map = glm_model.compute_contrast(contrast_val, output_type="z_score") # Use to produce and save only zscored effect size
        sub_map = glm_model.compute_contrast(contrast_val, output_type='all')
        sub_dict[sub+':'+contrast_id] = sub_map # Save to GLM result dictionary


#%% # Save the unthresholded data for all participants

''' Unthresholded zmaps will be required for group level GLM '''

# Call subject keys within the dictionary and save the zscored beta maps without thresholding
for key,data in sub_dict.items():    
    nib.save(sub_dict[key]['z_score'], f'/media/austin/Emotion/InterAffectTravel/neuralData/analysis/GLM/fstlvl/{key}_unfiltered_zmap.nii.gz') # change z_score to other data types to save them
    

#%% # pThresholding for all participant data

''' To save thresholded zscored effect sizes calculated by the first level GLM '''
# Loop across participants
for j in range(file_list.shape[0]):
    
    # Call subject keys within the dictionary and save the zscored beta maps without thresholding
    for key,data in sub_dict.items():
        
        # Check for participant number and call keys specific to the participant
        if f"{j}".zfill(2) in key:
            
            # Extract the zscored effect size map
            z_map = sub_dict[key]['z_score']

            # beta_map = sub_dict[key]['effect_size'] # Uncomment if you require non-zscored effect size values
            
            # Threshold the stat map(zscored effect size map) with 5% FDR and 10 voxel cluster threshold
            pThresh = nilearn.glm.threshold_stats_img(z_map, height_control = 'fdr', alpha = 0.05, cluster_threshold=10)[0]
            
            # Zero out voxels that are insignificant after FDR correction 
            FinImg = math_img("img1*img2", img1 = z_map, img2 = binarize_img(pThresh), copy_header_from = 'img1')
            
            #plot_stat_map(FinImg, cut_coords=[1,-40,33])
            
            # Save thresholded stat maps
            nib.save(FinImg, f'/media/austin/Emotion/InterAffectTravel/neuralData/analysis/GLM/zmap_AVNA_pThresh/{key}_pThresh_zmap.nii.gz')


#%% # Second level analysis, confirmation of group level behaviour in the participant pool

''' USE UNTHRESHOLED STAT MAPS FOR THIS. DO NOT USE FDR CORRECTED ONES AS THEY WILL ERROR OUT'''

''' 
    The second level GLM uses a design matrix with the subjects as the variable under consideration, making it a group level comparison.
    This allows for any group level effects to be extracted and observed.
'''

# glob to get location information for all unthresholded stat maps saved previously
file_list = sorted([x for x in glob.glob(os.path.join('/media/austin/Emotion/InterAffectTravel/neuralData/analysis/GLM/fstlvl/sub*A-NA*'))])

# Intialize empty array 
second_level_input = []

# Loop over participant files
for i,f in enumerate(file_list):    
    # Extract subject number
    sub = os.path.basename(f).split(':')[0]
    print('The following subject data is being processed:',sub)
    # Load subject specific uthresholded zmap
    contData = nib.load(f) 
    # Append to create a cluster of subject specific uthresholded zmap
    second_level_input.append(contData)

# Create second level design matrix 
design_matrix = pd.DataFrame( [1] * len(file_list), columns=["intercept"] )

# Initialize second level GLM 
glm2model = second_level.SecondLevelModel()
# Fit data to second level GLM 
glm2model = glm2model.fit(second_level_input, design_matrix = design_matrix)
# Compute group level effects
z_map = glm2model.compute_contrast(output_type='all')

# Threshold using a 1% FDR and 10 voxel cluster threshold. 
''' This mask produces the output for a one tailed hypothesis test. It can also produce two tailed results where any difference is shown by setting two_sided to false '''
pMask = nilearn.glm.threshold_stats_img(z_map['z_score'], height_control = 'fdr', alpha = 0.05, cluster_threshold=10, two_sided=False)[0]
# FinImg = math_img("img1*img2", img1 = z_map['z_score'], img2=binarize_img(pMask), copy_header_from='img1')
#pMask2 = nilearn.glm.threshold_stats_img(pMask, threshold=0, height_control = None)[0]
#plot_stat_map(z_map['z_score'], cut_coords=[1,-40,33])

plot_stat_map(pMask, cut_coords=[1,-40,33], title='Regions active in emotional tranistions')
plt.savefig("/home/austin/Documents/AustinV_Thesis/figures/finalFigure/GLM_TransitionCoding_A-NA.png", transparent=True)

nib.save(pMask, '/media/austin/Emotion/InterAffectTravel/neuralData/analysis/GLM/glm2ndLVL_A-NA_pThresh5%_zmap.nii.gz')

\end{lstlisting}

